\documentclass[12pt,a4paper]{article}
\usepackage[utf8]{inputenc}
\usepackage[indonesian]{babel}
\usepackage{amsmath}
\usepackage{amsfonts}
\usepackage{amssymb}
\usepackage{geometry}

\geometry{margin=2.5cm}

\title{No 14 - Single Answers Multiple Choice}
\author{Ahmad Fakhrul Bawani}
\date{}

\begin{document}

\maketitle

\section*{Soal A}

Diketahui fungsi $f(x)$ didefinisikan sebagai berikut:
\begin{equation*}
  f(x) = 
  \begin{cases} 
    \dfrac{\sin x}{x}, & \text{untuk } 0 < x \le \pi \\[1ex]
    1, & \text{untuk } x = 0 
  \end{cases}
\end{equation*}

Which of the following choices correctly shows that $x f(x) = \sin x$ for $0 \le x \le \pi$?

\subsection*{Pilihan Jawaban:}

\begin{itemize}
  \item[\textbf{A.}] For $0 \le x \le \pi$, the expression $\dfrac{\sin x}{x}$ is defined and so $x f(x) = x \left(\dfrac{\sin x}{x}\right) = \sin x$.
  
  \item[\textbf{B.}] For $0 < x \le \pi$, the expression $\dfrac{\sin x}{x}$ is defined and so $x f(x) = x \left(\dfrac{\sin x}{x}\right) = \sin x$. For $x = 0$, $x f(x) = x(1) = 0$. Since $\sin x = 0$ for $x = 0$, $x f(x) = \sin x$ for $0 \le x \le \pi$.
  
  \item[\textbf{C.}] For $0 < x \le \pi$, the expression $\dfrac{\sin x}{x}$ is defined and so $x f(x) = x \left(\dfrac{\sin x}{x}\right) = \sin x$. For $x = 0$, $x f(x) = x(1) = 1$. Since $\sin x = 1$ for $x = 0$, $x f(x) = \sin x$ for $0 \le x \le \pi$.
\end{itemize}

\subsection*{Kunci Jawaban Soal A}
Jawaban yang benar adalah \textbf{B}.

\section*{Soal B}

Find the volume of the solid generated by revolving the shaded region about the $y$-axis in the accompanying figure (dari $x = 0$ sampai $x = \pi$).

\subsection*{1. Rumus Umum (Metode Shell)}
Karena daerah di bawah kurva $y = f(x)$ pada interval $[0, \pi]$ diputar mengelilingi sumbu-$y$, volume benda putar menggunakan metode kulit tabung didefinisikan sebagai:
\begin{equation*}
  V = 2\pi \int_{a}^{b} (\text{jari-jari}) \cdot (\text{tinggi}) \, dx
\end{equation*}
Dimana:
\begin{itemize}
  \item Jari-jari kulit tabung adalah $x$.
  \item Tinggi kulit tabung adalah $f(x)$.
\end{itemize}

\subsection*{2. Solusi Integrasi}
Substitusikan komponen ke dalam rumus:
\begin{equation*}
  V = 2\pi \int_{0}^{\pi} x \cdot f(x) \, dx
\end{equation*}

Berdasarkan pembuktian pada Soal A, kita tahu bahwa nilai $x f(x) = \sin x$ untuk seluruh interval $[0, \pi]$. Oleh karena itu, integralnya dapat disederhanakan menjadi:
\begin{align*}
  V &= 2\pi \int_{0}^{\pi} \sin x \, dx \\
  V &= 2\pi \Big[ -\cos x \Big]_{0}^{\pi} \\
  V &= 2\pi \Big( -\cos(\pi) - (-\cos(0)) \Big) \\
  V &= 2\pi \Big( -(-1) + 1 \Big) \\
  V &= 2\pi (1 + 1) \\
  V &= 2\pi (2) \\
  V &= 4\pi
\end{align*}

Jadi, volume benda putar yang dihasilkan adalah \textbf{$4\pi$ satuan volume}.

\end{document}